% !PRINTED 618
\weberpage{618}
%<<CONT
is to be set. $r$ is therefore an \gesperrt{integer} in $\Omega_m$, and its norm is $=+1$; consequently $r$ is a unit.

$\Omega_m$ is a normal field, which under the $\mu$ substitutions
\[ s_n = (r, r^n) \]
passes over into itself.

These $\mu$ substitutions $s_n$ form an \gesperrt{Abelian} group $\mathfrak{N}$, which is the \gesperrt{Galois} group of the field $\Omega_m$. It is isomorphic with the group, likewise to be denoted by $\mathfrak{N}$, of the number classes $n$ taken with respect to the modulus $m$ (\S.\,16).

The function $f(x)$ can be decomposed into the $\mu$ factors $x - r^n$, and we therefore set
\weq{4}{f(x) = \prod\limits^{n} (x - r^n).}

If we set $x=1$ in this, and observe that [by (3)] $f(1) = q$, it follows that
\weq{5}{q = \prod\limits^{n} (1 - r^n).}

The numbers
\weq{6}{\sigma_n = 1 - r^n, \quad \sigma = 1 - r}
are, however, integers of the field $\Omega_m$, and formula (5) shows first of all \gesperrt{that $q$ is decomposable in the field $\Omega_m$ into $\mu$ factors $\sigma_n$.}

We prove first that the $\mu$ numbers $\sigma_n$ are associated with one another. If $n, n'$ denote two numbers not divisible by $q$, we can determine a natural number $a$ in such a way that
\[ n' \equiv a n \ (\mathrm{mod}\ m) \]
holds. But then
\[ \frac{\sigma_{n'}}{\sigma_n} = \frac{1 - r^{an}}{1 - r^n} = 1 + r^n + r^{2n} + \cdots + r^{(a-1)n} \]
is an integer, and hence $\sigma_{n'}$ is divisible by $\sigma_n$. Since $n$ can here be interchanged with $n'$, $\sigma_n$ is also divisible by $\sigma_{n'}$, and hence the two numbers are associated (\S.\,138). If therefore $\varepsilon$ denotes a unit, then by (5):
\weq{7}{q = \varepsilon \sigma^{\mu}, \quad N(\sigma) = q.}

Thus the natural prime $q$ is associated with the $\mu$th power of an integer $\sigma$ in the field $\Omega_m$. But this number $\sigma$ is still a prime in the field $\Omega_m$ as well, and indeed of the first degree. For if $\sigma$ has any divisor $\sigma'$, the norm of $\sigma'$ must be a divisor of the norm of $\sigma$; hence, if $\sigma'$ is not a
%CONT>>
